\documentclass[10pt,a4paper]{amsart}
\usepackage[margin=19mm]{geometry}
\usepackage{amsmath,amssymb,mathtools}
\usepackage[scr=rsfs,cal=euler]{mathalfa}
\usepackage{xfrac}
\usepackage[unicode,hidelinks,bookmarks=false]{hyperref}
\newcommand{\ie}{i.e.\ }
\newcommand{\cf}{cf.\ }
\newcommand{\resp}{respectively\ }
\newcommand{\Subsection}{\subsection}
\providecommand{\nobreakdash}{\nobreak\hbox{-}}
\setlength{\emergencystretch}{2em}
\multlinegap0pt
\usepackage[shortlabels]{enumitem}
\usepackage{amssymb}
\usepackage{mathtools}
\usepackage{bbm}
\usepackage{tikz}
\usepackage{xcolor}
\usepackage[normalem]{ulem}
\newtheorem{lemma}{Lemma}
\newcommand{\F}{\mathbb{F}}
\newcommand{\e}{\mathbbm{1}}          % indicator function
\title{\vspace{-0.7cm} Roth-type theorems in $K_{s,t}$-free sets}
\author{Yifan Jing, Cosmin Pohoata, Max Wenqiang Xu}
\date{Source: arXiv:2601.18738v1 (CC BY 4.0)}
\begin{document}
\maketitle

\noindent\emph{Source and licence.} \vspace{-0.7cm} Roth-type theorems in $K_{s,t}$-free sets. Version arXiv:2601.18738v1 (CC BY 4.0), distributed by arXiv under the Creative Commons Attribution 4.0 International licence, http://creativecommons.org/licenses/by/4.0/, as recorded in the arXiv licence metadata for this exact version and shown on the abstract page https://arxiv.org/abs/2601.18738v1. Full licence notice: \texttt{licenses/arxiv-2601.18738-CC-BY-4.0.txt}.\par\smallskip

\noindent\textit{Editorial scope.} Counted unit: the article's lemma lem:dense-model-ff (source0.tex lines 863--871). The article's complete original proof of it is delivered with it (source0.tex lines 873--1009, one \texttt{proof} environment of 137 lines). No mathematics is written, repaired or re-derived: the statement and the proof are the article's own lines, in the article's own notation, and no retained line is substituted or re-flowed.  Retained uncounted as the source-local context the proof needs: lines 290--297; lines 443--448; lines 481--486.  Nothing was omitted from any retained span, so the package carries no omission record.  No retained line contains a citation, so the wrapper carries no bibliography and no bibliography entry.  No retained line contains a figure environment, an image-inclusion command or drawing code, so no image is delivered and the package declares no asset dependency.  Editorial formatting only, in addition: an AMS \texttt{amsart} preamble that restates the article's own declarations and macros used by the retained lines, namely the environments \texttt{lemma} on separate counters, together with the standard packages those lines need.  The retained environments print in this document as Lemma 1 in order of appearance, whereas the article prints the same items with the numbering of its own full document.  The complete native source of the article as distributed in the arXiv e-print for this exact version is bundled unchanged as \texttt{sources/193413/source0.tex}.
\begin{lemma}\label{lem: E2}
Suppose $A\subseteq [N]$ satisfies $E_s(\e_A,\e_{-A})\ll |A|^s$ for integer  $s\ge 2$. Then
\[
E_{2}(\e_A,\e_{-A})\ll |A|^{3-\frac{1}{s-1}},
\qquad
E_{s-1}(\e_A,\e_{-A})\ll |A|^{s-\frac{s-2}{s-1}}.
\]
\end{lemma}

\begin{lemma}\label{lem: upper bound on A} Let $\eta\in(0,1)$  and $2\le s\le t$.
Suppose $A\subseteq [N]$ and $E_s(\e_A,\e_{-A})\leq (t+\eta)|A|^s$. Then
\[
|A|\leq 2\Bigl(\frac{t^2}{1-\eta}\Bigr)^{1/s}N^{1-1/s}.
\]
\end{lemma}

\begin{lemma}\label{lem: vanishing K_st-free}
Let $A\subseteq[N]$ be $K_{s,t}$-free. Then\[
\sum_{a_1,\dots,a_s\in A} (r_A(a_1,\dots,a_s) - (t-1))_+\ll_{s,t}|A|^{s-c_s},
\]
for some constant $c_s\in(0,1)$ depending only on $s$. 
\end{lemma}

\begin{lemma}[Dense model in $\F_q^n$]\label{lem:dense-model-ff}
Suppose $A\subseteq G$ satisfies $|A|>\delta N^{1-1/s}$ and $A$ is $K_{s,t}$-free. 
Then for every $0<\varepsilon<1$ there exists $f:G\to[0,\infty)$ such that
\begin{enumerate}[label=(\roman*)]
\item $\sum_{x\in G} f(x)=N^{1/s}|A|$;
\item $\|\widehat f-N^{1/s}\widehat{\e_A}\|_\infty\ll \varepsilon N$;
\item $\displaystyle \sum_{x\in G} f(x)^s\le N\Big(t+\exp(\varepsilon^{-O_s(1)})\cdot O_{s,t}(\eta)\Big)$ for some $\eta\ll_{s,t}|A|^{-c_s}$.
\end{enumerate}
\end{lemma}

\begin{proof}
Define the large spectrum
\[
\operatorname{Spec}(A,\varepsilon):=\{\xi\in G:\ |\widehat{\e_A}(\xi)|\ge \varepsilon |A|\}.
\]
Let $V:=\mathrm{Span}(\operatorname{Spec}(A,\varepsilon))\le G$ be the $\F_q$-linear span of $\operatorname{Spec}(A,\varepsilon)$, and set
\[
H:=V^\perp=\{x\in G:\ \chi_\xi(x)=1 \text{ for all }\xi\in V\}.
\]
Then $H$ is a subspace of $G$ and $\widehat{\mu_H}=\e_V$ (i.e.\ $\widehat{\mu_H}(\xi)=1$ if
$\xi\in V$ and $0$ otherwise).

We next claim that $\dim(V)\ll_{s,t,\delta}\varepsilon^{-4}$. To see this, let $\Lambda\subseteq \operatorname{Spec}(A,\varepsilon)$ be any set of linearly independent vectors spanning
$V$, so $\dim(V)=|\Lambda|$.  Then
\[
|\Lambda|\,\varepsilon^4|A|^4\ \le\ \sum_{\xi\in\Lambda}|\widehat{\e_A}(\xi)|^4
\ \le\ \sum_{\xi\in G}|\widehat{\e_A}(\xi)|^4
\ =\ N\cdot E_2(\e_A,\e_{-A}).
\]
By Lemma~\ref{lem: E2} and the hypothesis
$E_s(\e_A,\e_{-A})\ll_{s,t}|A|^s$, we have
$E_2(\e_A,\e_{-A})\ll_s |A|^{3-\frac{1}{s-1}}$, so
\[
\dim(V) = |\Lambda| \ \ll_s\ \frac{N}{\varepsilon^4|A|^{1+\frac{1}{s-1}}}.
\]
Using $|A|>\delta N^{1-1/s}$ and the identity
$(1-\tfrac1s)(1+\tfrac{1}{s-1})=1$, we obtain $\dim(V)\ll_{s,\delta}\varepsilon^{-4}$.
In particular,
\begin{equation}\label{eqn: H}
|H|=q^{n-\dim(V)}\ge q^{-O_{s,\delta}(\varepsilon^{-4})}N=\exp(-\varepsilon^{-O_s(1)})N.    
\end{equation}
%(absorbing $q$ and $\delta$ into the implied constants).

%\subsection*{The function $f$} 
Let $\mu_H:=|H|^{-1}\e_H$ and define
\[
f:=N^{1/s}\,\e_A*\mu_H.
\]
We now check that $f$ satisfies properties (i)-(iii). Clearly, $f\ge0$ and (i) is immediate since $\sum_x \mu_H(x)=1$. For (ii), using $\widehat{\mu_H}=\e_V$ we have for every $\xi\in G$,
\[
\widehat f(\xi)=N^{1/s}\widehat{\e_A}(\xi)\widehat{\mu_H}(\xi)
=
\begin{cases}
N^{1/s}\widehat{\e_A}(\xi),& \xi\in V,\\
0,& \xi\notin V.
\end{cases}
\]
Hence $\widehat f(\xi)-N^{1/s}\widehat{\e_A}(\xi)=0$ for $\xi\in V$, while for $\xi\notin V$ we have
$\xi\notin \operatorname{Spec}(A,\varepsilon)$ and therefore $|\widehat{\e_A}(\xi)|<\varepsilon|A|$.  Thus
\[
\|\widehat f-N^{1/s}\widehat{\e_A}\|_\infty
\le N^{1/s}\varepsilon|A|
\ll_{s,t}\varepsilon N,
\]
where in the last step we used a finite field version of  Lemma~\ref{lem: upper bound on A} to bound $|A|\ll_{s,t}N^{1-1/s}$.

To prove the $L^s$ bound (iii), we start by writing $\e_H$ for the (unnormalized) indicator of $H$ and by noting that
$f^s=N(\e_A*\mu_H)^s$.  Therefore
\begin{equation}\label{eqn: fs}
\sum_{x\in G} f(x)^s
=
N\sum_{x\in G} (\e_A*\mu_H(x))^s
=
N|H|^{-s}\sum_{x\in G} (\e_A*\e_H(x))^s.  
\end{equation}
For convenience, set
\[
S:=\sum_{x\in G} (\e_A*\e_H(x))^s,
\]
and rewrite the quantity as follows:
\[
S=\sum_{h_1,\dots,h_s\in H}\ \sum_{x\in G}\ \prod_{i=1}^s \e_A(x-h_i)
=\sum_{h_1,\dots,h_s\in H} |S_{\mathbf{h}}|,
\]
where $S_{\mathbf{h}}:=\{x\in G:\ x-h_i\in A\,\ \text{for all}\ 1\le i \le s\}$ and $\mathbf{h}=(h_1,\dots,h_s)$.

For an ordered $s$-tuple $\mathbf{a}=(a_1,\dots,a_s)\in A^s$ define
\[
r_A(\mathbf{a})
:=
\bigl|\{d\in G\setminus\{0\}:\ a_i-d\in A \text{ for all }1\le i\le s\}\bigr|.
\]
As $A$ is $K_{s,t}$-free, Lemma~\ref{lem: vanishing K_st-free} implies that there is $\eta = O(|A|^{-c_s})$ such that 
\begin{equation}\label{eq:excess_rA}
\sum_{\mathbf{a}\in A^s}\bigl(r_A(\mathbf{a})-(t-1)\bigr)_+ \ \le\ \eta |A|^s,
\end{equation}
 where $(x)_+ := \max\{x,0\}$ for a real number $x$. 

Next, fix $\mathbf{h}$ and suppose $S_{\mathbf{h}}\neq\emptyset$.  For any $x\in S_{\mathbf{h}}$ the
associated tuple $\mathbf{a}(x):=(x-h_1,\dots,x-h_s)\in A^s$ satisfies
\begin{equation}\label{eq:rA_equals_size}
r_A(\mathbf{a}(x))=|S_{\mathbf{h}}|-1.
\end{equation}
Indeed, if $x'\in S_{\mathbf{h}}$ then $\mathbf{a}(x')=\mathbf{a}(x)-(x-x')$ coordinatewise, giving a shift
counted by $r_A$; conversely, if $\mathbf{a}(x)-d\in A^s$ then $x-d\in S_{\mathbf{h}}$.
In particular, $|S_{\mathbf{h}}|>t$ if and only if $r_A(\mathbf{a}(x))>t-1$.

Notice that $|S_{\mathbf{h}}|\ \le\ t+\bigl(|S_{\mathbf{h}}|-t\bigr)_+$ always holds and it implies that
\begin{equation}\label{eqn: S}
   S\ \le\ t|H|^s+\sum_{\mathbf{h}\in H^s}\bigl(|S_{\mathbf{h}}|-t\bigr)_+. 
\end{equation}
If $|S_{\mathbf{h}}|>t$, then $|S_{\mathbf{h}}|\bigl(|S_{\mathbf{h}}|-t\bigr)_+\ge
t\bigl(|S_{\mathbf{h}}|-t\bigr)_+$, and further
\begin{equation}\label{eq:excess_linear}
\sum_{\mathbf{h}\in H^s}\bigl(|S_{\mathbf{h}}|-t\bigr)_+
\ \le\ \frac1t\sum_{\mathbf{h}\in H^s}|S_{\mathbf{h}}|\bigl(|S_{\mathbf{h}}|-t\bigr)_+.
\end{equation}
Using \eqref{eq:rA_equals_size}, the right hand side of \eqref{eq:excess_linear} equals
\[
\frac1t\sum_{\mathbf{h}\in H^s}\ \sum_{x\in S_{\mathbf{h}}}\bigl(r_A(\mathbf{a}(x))-(t-1)\bigr)_+.
\]
Group this sum by $\mathbf{a}\in A^s$.  For a fixed $\mathbf{a}=(a_1,\dots,a_s)$, the number of pairs
$(\mathbf{h},x)$ with $x-h_i=a_i$ for all $i$ is at most $|H|$: indeed, choosing any $h_1\in H$ determines
$x=a_1+h_1$ and then $h_i=x-a_i=h_1+(a_1-a_i)$. Therefore
\[
\sum_{\mathbf{h}\in H^s}|S_{\mathbf{h}}|\bigl(|S_{\mathbf{h}}|-t\bigr)_+
\ \le\ |H|\sum_{\mathbf{a}\in A^s}\bigl(r_A(\mathbf{a})-(t-1)\bigr)_+
\ \le\ \eta |A|^s |H|
\]
where in the last inequality we used
\eqref{eq:excess_rA}.  Combining this with \eqref{eq:excess_linear} and \eqref{eqn: S} gives
\[
S\ \le\ t|H|^s+\frac{\eta}{t}|A|^s|H|.
\]
Invoking \eqref{eqn: fs} gives
\[
\sum_{x\in G} f(x)^s
\le
N\left(t+\frac{\eta}{t}\cdot\frac{|A|^s}{|H|^{s-1}}\right).
\]
Finally, using $|A|\ll_{s,t}N^{1-1/s}$ from Lemma~\ref{lem: upper bound on A} and
$|H|\ge \exp(-\varepsilon^{-O_s(1)})N$ from \eqref{eqn: H} , we conclude that
$
\frac{|A|^s}{|H|^{s-1}}
\ \ll_{s,t}\ \exp(\varepsilon^{-O_s(1)})$,
and thereby (iii) follows.
\end{proof}

\end{document}
